\documentclass[10pt]{article}
\usepackage[top=1.1in, bottom=1.1in, left=1.1in, right=1.1in, letterpaper]{geometry}
\usepackage{algorithm}
\usepackage{algorithmic}
\usepackage[numbers,compress]{natbib}
% 共享文档信息 (可以从主文档和补充文档中调用)

% Encoding and typography (load early)
\usepackage[utf8]{inputenc} % UTF-8 input encoding
\usepackage[T1]{fontenc}    % 8-bit T1 fonts
\usepackage{microtype}      % Microtypography improvements

% General utilities
\usepackage{lipsum}         % Placeholder text
\usepackage{graphicx}       % Image insertion
% \usepackage{epstopdf}     % Uncomment only if EPS conversion is explicitly needed
\usepackage{xcolor}         % Colors
\usepackage{url}            % URL typesetting
\usepackage{nicefrac}       % Compact fraction symbols
\usepackage{xspace}         % Smart spacing for macros

% Mathematical packages
\usepackage{amsmath}        % Core math support
% \theoremstyle{remark}
\usepackage{amssymb}        % Math symbols (includes amsfonts)
\usepackage{mathrsfs}       % Script fonts
\usepackage{bm}             % Bold math symbols
\usepackage{booktabs}       % Professional tables
\usepackage{subcaption}     % Subfigures
\usepackage{multirow}       % Multirow tables
\usepackage{array}          % Enhanced table formatting
\usepackage{setspace}       % Line spacing control

% Algorithms
\usepackage{algorithm}
\usepackage{algorithmic}    % Consider algorithmicx or algorithm2e for modern alternatives
\usepackage{tikz}
\usetikzlibrary{arrows.meta,positioning,fit,shapes.misc}
% References and hyperlinks (load late)
\usepackage{hyperref}       % Hyperlinks
\usepackage{nameref}        % Reference names
\usepackage{cleveref}       % Smart cross-references (loaded after amsmath and hyperref)

% Ensure hyperref settings are robust
\hypersetup{
    colorlinks=true,
    linkcolor=blue,
    citecolor=blue,
    urlcolor=blue
}
%-----------------------------数学符号新命令------------------------------
\newcommand{\eucalB}{\EuScript{B}}
\newcommand{\eucalP}{\EuScript{P}}
\newcommand{\Var}{\mathrm{Var}}
\newcommand{\Cov}{\mathrm{Cov}}
\newcommand{\Ex}{\mathsf{E}}
\newcommand{\Prob}{\mathsf{P}}
\newcommand{\one}{\mathbbm{1}}
\newcommand{\One}{\bm{1}}
\newcommand{\zero}{\bm{0}}
\newcommand{\RR}{\mathbb{R}}
\newcommand{\ZZ}{\mathbb{Z}}
\newcommand{\EE}{\mathbb{E}}
\newcommand{\PP}{\mathbb{P}}
\newcommand{\Rp}{\RR_+}
\newcommand{\TT}{\mathbb{T}}
\newcommand{\NN}{\mathbb{N}}
\newcommand{\VV}{\mathbb{V}}
\newcommand{\diff}{\mathrm{d}}
\newcommand{\prox}{\mathrm{prox}}
\newcommand{\rprox}{\mathrm{rprox}}
\newcommand{\ran}{\mathrm{ran}}
\newcommand{\Id}{\mathrm{Id}}
\newcommand{\e}{\mathrm{e}}
\newcommand{\bA}{\bm{A}}
\newcommand{\bB}{\bm{B}}
\newcommand{\bH}{\bm{H}}
\newcommand{\bF}{\bm{F}}
\newcommand{\bS}{\bm{S}}
\newcommand{\bX}{\bm{X}}
\newcommand{\bx}{\bm{x}}
\newcommand{\bb}{\bm{b}}
\newcommand{\bc}{\bm{c}}
\newcommand{\br}{\bm{r}}
\newcommand{\bbeta}{\bm{\beta}}
\newcommand{\bt}{\bm{t}}
% \newcommand{\eproof}{\hfill $\square$}
\newcommand{\bh}{\bm{h}}
\newcommand{\bu}{\bm{u}}
\newcommand{\bz}{\bm{z}}
\newcommand{\ba}{\bm{a}}
\newcommand{\by}{\bm{y}}
\newcommand{\be}{\bm{e}}
\newcommand{\bU}{\bm{U}}
\newcommand{\bZ}{\bm{Z}}
\newcommand{\bY}{\bm{Y}}
\newcommand{\blmbd}{\bm{\lambda}}
\newcommand{\bLmbd}{\bm{\Lambda}}
\newcommand{\bg}{\bm{g}}
\newcommand{\bd}{\bm{d}}
\newcommand{\bmu}{\bm{\mu}}
\DeclareMathOperator*{\argmin}{\mathrm{argmin}}
\DeclareMathOperator*{\Argmin}{\mathrm{Argmin}}
\DeclareMathOperator*{\minimize}{\mathrm{minimize}\,}
\DeclareMathOperator*{\maximize}{\mathrm{maximize}}
\DeclareMathOperator*{\argmax}{\mathrm{argmax}}
\DeclareMathOperator*{\Diag}{\mathrm{Diag}}
\newcommand{\calX}{\mathcal{X}}
\newcommand{\calY}{\mathcal{Y}}
\newcommand{\calB}{\mathcal{B}}
\newcommand{\calL}{\mathcal{L}}
\newcommand{\calG}{\mathcal{G}}
\newcommand{\calS}{\mathcal{S}}
\newcommand{\calC}{\mathcal{C}}
\newcommand{\calJ}{\mathcal{J}}
\newcommand{\calD}{\mathcal{D}}
\newcommand{\calQ}{\mathcal{Q}}
\newcommand{\calO}{\mathcal{O}}
\newcommand{\ent}{\mathrm{ent}}
% \newcommand{\diag}{\mathrm{diag}}
\newcommand{\dom}{\mathrm{dom}}
\newcommand{\interior}{\mathrm{int}}
\newcommand{\zer}{\mathrm{zer}}
\newcommand{\epi}{\mathrm{epi}}
\newcommand{\amax}{a_{\max}}
\newcommand{\amin}{a_{\min}}
\newcommand{\bmax}{b_{\max}}
\newcommand{\bmin}{b_{\min}}
\newcommand{\bp}{\bm{p}}
\newcommand{\bD}{\bm{D}}
\newcommand{\mW}{\mathrm{W}}
\newcommand{\bv}{\bm{v}}
\newcommand{\bw}{\bm{w}}
\newcommand{\ta}{\tilde{a}}
\newcommand{\tb}{\tilde{b}}
\newcommand{\tp}{\widetilde{\bp}}
\newcommand{\sumn}{\sum_{i=1}^n}
\newcommand{\sumQ}{\sum_{q=1}^Q}
\newcommand{\sumq}{\sum_{k=1}^q}
\newcommand{\summ}{\sum_{i=1}^m}
\newcommand{\sumd}{\sum_{i=1}^d}
\newcommand{\sumT}{\sum\limits_{ k =1}^{T}}
\newcommand{\sumTa}{\frac{1}{T} \sum\limits_{ k =1}^{T}}
\newcommand{\bC}{\bm{C}}
\newcommand{\bI}{\bm{I}}
\newcommand{\bM}{\bm{M}}
\newcommand{\bR}{\bm{R}}
\newcommand{\bV}{\bm{V}}
\newcommand{\bW}{\bm{W}}
\newcommand{\bxi}{\bm{\xi}}
\newcommand{\tmu}{\widetilde{\bmu}}
\newcommand{\tU}{\widetilde{\bU}}
\newcommand{\tV}{\widetilde{\bV}}
\newcommand{\tM}{\widetilde{\bM}}
\newcommand{\bDelta}{\bm{\Delta}}
\newcommand{\bs}{\bm{s}}
\newcommand{\ox}{\overline{x}}
\newcommand{\ob}{\overline{b}}
\newcommand{\op}{\overline{p}}
\newcommand{\obp}{\overline{\bp}}
\newcommand{\obmu}{\overline{\bmu}}
\newcommand{\obC}{\overline{\bC}}
\newcommand{\obD}{\overline{\bD}}
\newcommand{\bsigma}{\bm{\sigma}}
\newcommand{\bomega}{\bm{\omega}}
\newcommand{\tomega}{\widetilde{\omega}}
\newcommand{\tr}{\mathrm{tr}}
\newcommand{\an}{\mathrm{an}}
\newcommand{\bj}{\bm{j}}
\newcommand{\bG}{\bm{G}}
\newcommand{\bcalG}{\bm{\calG}}
\newcommand{\bcalK}{\bm{\calK}}
\newcommand{\bcalX}{\bm{\calX}}
\newcommand{\fX}{\mathfrak{X}}
\newcommand{\bfX}{\bm{\fX}}
\newcommand{\bell}{\bm{\ell}}
\newcommand{\bL}{\bm{L}}
\newcommand{\bzeta}{\bm{\zeta}}
\renewcommand{\vec}{\mathrm{vec}}
\newcommand{\hy}{\widehat{y}}
\newcommand{\ul}[1]{\underline{#1}}
\newtheorem{assumption}{Assumption}
\newtheorem{definition}{Definition}
\newtheorem{lemma}{Lemma}
\newtheorem{theorem}{Theorem}
\newtheorem{proof}{Proof}
\newenvironment{noproof}{\par\noindent\textbf{Proof.}\space}{\hfill\par}
\newtheorem{corollary}[theorem]{Corollary} 
\newtheorem{remark}{Remark}[section]  % 按节编号的注释环境
\newtheorem{proposition}{Proposition}

% \newtheorem{remark}[theorem]{Remark}  % 与theorem共享编号

\newcommand{\bcL}{\overline{\calL}}
\newcommand{\KL}{K\L\xspace}
\newcommand{\res}{\mathrm{res}}
\newcommand{\eproof}{\hfill $\square$}
\newcommand{\calF}{\mathcal{F}}

\newcommand{\commyy}[1]{{\color{red}(yy: #1)}} % Yuan Yao's comments

\numberwithin{equation}{section}
\newcommand{\blue}[1]{\begin{color}{blue}#1\end{color}} 
\newcommand{\red}[1]{\begin{color}{red}#1\end{color}}

%% Packages and macros go here
\usepackage{lipsum}
\usepackage{amsfonts}
\usepackage{graphicx}
\usepackage{epstopdf}
\usepackage{algorithmic}
\ifpdf
  \DeclareGraphicsExtensions{.eps,.pdf,.png,.jpg}
\else
  \DeclareGraphicsExtensions{.eps}
\fi

% Prevent itemized lists from running into the left margin inside theorems and proofs
\usepackage{enumitem}
\setlist[enumerate]{leftmargin=.5in}
\setlist[itemize]{leftmargin=.5in}

% Add a serial/Oxford comma by default.
\newcommand{\creflastconjunction}{, and~}


\usepackage{multibib}
\newcites{supp}{References (Supplementary Material)}

\usepackage{wrapfig}
\usepackage{caption}
% \usepackage{subcaption}
\usepackage[title,toc,page]{appendix}
\pdfstringdefDisableCommands{%
  \def\Cref#1{Ref #1}%
  \let\new@ifnextchar\relax%
  \let\textsubscript\relax% 处理下标
  \let\textsuperscript\relax% 处理上标
  \def\_{}% 处理下划线
}
\usepackage{multirow}
\allowdisplaybreaks[2]  % 中等优先级允许公式跨页
\usepackage{xcolor} % 在导言区加载

% \captionsetup[figure]{compatibility=false}

% \PassOptionsToPackage{numbers, sort}{natbib} % 参考文献 
% \PassOptionsToPackage{numbers, sort, square}{natbib}
\title{Stochastic Perturbed Proximal Augmented Lagrangian for Nonconvex Constrained Optimization}
\newcommand{\calE}{\mathcal{E}}
\newcommand{\calI}{\mathcal{I}}

\begin{document}

\maketitle

\section{Introduction}
\label{sec:intro}
We consider the following nonconvex stochastic optimization problem with
mixed deterministic constraints:
\begin{align}\label{Prob: primal problem}
    \min_{\bx \in \calX} \quad & F(\bx) := f(\bx) + g(\bx), \notag \\
    \text{s.t.} \quad & \bc_{\calE}(\bx) = \bm{0}, \notag\\
    & \bc_{\calI}(\bx) \le \bm{0}, \tag{{\bf P}}
\end{align}
where $\calE$ and $\calI$ index the equality and inequality blocks,
respectively, and
$\bc(\bx):=\big(\bc_{\calE}(\bx),\bc_{\calI}(\bx)\big)\in\RR^m$
with $m:=|\calE|+|\calI|$.
Here $f(\bx)=\EE_{\xi}[\mathbf{f}(\bx,\xi)]$ is a smooth nonconvex
expectation objective, $g:\RR^n\to\RR\cup\{+\infty\}$ is a proper closed
proximal-friendly function, and the mappings
$\bc_{\calE}:\RR^n\to\RR^{|\calE|}$ and
$\bc_{\calI}:\RR^n\to\RR^{|\calI|}$ are continuously differentiable.
The set $\calX\subseteq\RR^n$ is nonempty and closed, the feasible set of
\eqref{Prob: primal problem} is nonempty, and $F$ is bounded below by
$F^*>-\infty$. Exact evaluation of $\nabla f$ is often impractical, so we
assume access only to stochastic gradients queried at sample points.

A key obstacle in this setting is that equality and inequality constraints
enter augmented-Lagrangian models asymmetrically: equality blocks generate the
standard quadratic augmentation, whereas inequality blocks require one-sided
activation. Our approach is to absorb this asymmetry into a single
\emph{effective multiplier}, which leads to a compact perturbed proximal
augmented Lagrangian with a fixed penalty parameter and a genuinely
single-loop primal--dual update.

\subsection{Related Work}
\label{sec:related-work}

\paragraph{Deterministic nonconvex constrained optimization.}
For deterministic nonconvex problems with equality and inequality constraints,
first-order methods with non-asymptotic complexity guarantees have been
studied extensively.
Proximal point methods~\cite{boob2023stochastic,ma2022quadratically} achieve
$\calO(\varepsilon^{-3})$ iteration complexity for smooth objectives via
strongly convex subproblems, and Boob et al.~\cite{boob2022level} further
attain $\calO(\varepsilon^{-2})$ by constructing quadratic models with
adaptive constraint levels.
Among penalty and augmented Lagrangian methods (ALMs), most approaches require
the penalty parameter to grow as $\Omega(\varepsilon^{-1})$, yielding
$\calO(\varepsilon^{-3})$ iteration complexity~\cite{li2021rate,jia2023first}.
Xie and Wright~\cite{xie2021proximal} achieve $\calO(\varepsilon^{-2})$
outer-loop complexity with a proximal ALM under a constant merit parameter,
but the algorithm is double-loop and the merit parameter must satisfy
$\rho = \Omega(\varepsilon^{-\eta})$ for some $\eta \in (0,1)$, making it
$\varepsilon$-dependent in theory.
SQP methods~\cite{curtis2021sqp,cartis2014complexity} also attain
$\calO(\varepsilon^{-2})$ complexity.

\paragraph{Stochastic nonconvex constrained optimization.}
We organize the literature by the two canonical settings:
\emph{semi-stochastic} (stochastic objective, deterministic constraints) and
\emph{fully stochastic} (both objective and constraints involve stochastic
oracles).

\emph{Semi-stochastic problems.}
Without additional structure on the noise, proximal
point~\cite{boob2023stochastic,ma2022quadratically},
penalty~\cite{jin2022nested}, and SQP~\cite{curtis2021sqp} methods all achieve
$\calO(\varepsilon^{-4})$ oracle complexity in stochastic gradient
evaluations, matching the $\calO(\varepsilon^{-4})$ lower bound for
unconstrained stochastic optimization~\cite{arjevani2023lower}.
Under a mean-squared smoothness (MSS) condition, variance reduction
techniques---STORM~\cite{cutkosky2019storm} and
SPIDER~\cite{fang2018spider}---improve the complexity to
$\calO(\varepsilon^{-3})$.
Representative results include the variance-reduced linearized ALM of
Shi et al.~\cite{shi2022linearized}, the truncated variance reduction approach
of Lu et al.~\cite{lu2022single}, and the inexact ALM of
Li et al.~\cite{li2021rate}, each attaining $\calO(\varepsilon^{-3})$ oracle
complexity under MSS for nonlinear constraints.

\emph{Fully stochastic problems.}
When constraint evaluations are also stochastic, complexity results are
uniformly higher.
Proximal point methods~\cite{boob2023stochastic,ma2022quadratically} incur
$\calO(\varepsilon^{-6})$ total oracle complexity due to the cost of solving
stochastic strongly convex subproblems.
Penalty-based ALMs employing STORM-type variance
reduction~\cite{alacaoglu2023augmented,cui2022complexity} face a large merit
parameter $\beta = \Theta(\varepsilon^{-1})$, which amplifies the variance of
the stochastic gradient and yields $\calO(\varepsilon^{-5})$ complexity even
under MSS.
Exact penalty methods combined with SPIDER~\cite{liu2022exact} achieve
$\calO(\varepsilon^{-4})$ (sub)gradient and $\calO(\varepsilon^{-6})$
constraint-evaluation complexity for nonsmooth problems.
Cui et al.~\cite{cui2024adaptive} recently propose an adaptive exact penalty
method for smooth \emph{equality}-constrained problems, obtaining
$\calO(\varepsilon^{-3})$ gradient and $\calO(\varepsilon^{-5})$
constraint-evaluation complexity; however, the analysis is confined to
equality constraints and relies on an exact penalty reformulation.
Stochastic SQP for equality-constrained fully stochastic
problems~\cite{na2022sqp} currently achieves only $\calO(\varepsilon^{-8})$.

\paragraph{ALMs with fixed penalty parameter.}
A fundamental challenge in ALM design for nonlinear constraints is that
standard convergence analysis requires the penalty parameter to scale as
$\Omega(\varepsilon^{-1})$.
This growth inflates the Lipschitz constant of the augmented Lagrangian and
forces the step size to shrink proportionally, degrading the stochastic oracle
complexity.
For polyhedral (linear) constraints, the Luo--Pang global error
bound~\cite{luo1994error} decouples feasibility control from the penalty
magnitude.
Exploiting this structure, Zhang and Luo~\cite{zhang2022proximal} achieve
$\calO(\varepsilon^{-2})$ iteration complexity for deterministic nonconvex
problems under a genuinely constant penalty.
In the stochastic setting, Alacaoglu and
Wright~\cite{alacaoglu2024complexity} combine a fixed penalty with variance
reduction to attain the optimal $\calO(\varepsilon^{-3})$ sample complexity
for linear equality constraints.
Huang et al.~\cite{huang2025stochastic} extend the fixed-penalty result to
linear inequality constraints via Moreau envelope smoothing, also reaching
$\calO(\varepsilon^{-3})$.
For general nonlinear constraints, existing stochastic methods either require
an increasing penalty parameter~\cite{li2024stochastic,alacaoglu2024complexity},
yielding $\calO(\varepsilon^{-4})$ oracle complexity, or introduce dual
momentum to partially compensate, achieving
$\calO(\varepsilon^{-3.5})$~\cite{shi2025momentum}.
To date, no single-loop ALM with a genuinely constant penalty parameter has
been shown to achieve $\calO(\varepsilon^{-3})$ oracle complexity for general
nonlinear mixed constraints in the stochastic setting.

Our perturbed ALM addresses this gap by combining a uniform proximal
perturbation on the multiplier vector with an effective-multiplier map that
absorbs the equality/inequality asymmetry. This yields a compact PPAL model,
keeps the penalty parameter fixed, and supports a single-loop stochastic
scheme for nonlinear mixed constraints.

\subsection{Our Contributions}
\begin{itemize}
    \item \textbf{Unified PPAL formulation for mixed constraints.}
    We introduce an effective multiplier
    $\hat{\blmbd}(\blmbd,\bx)$ that absorbs the difference between equality
    and inequality blocks, allowing the perturbed proximal augmented
    Lagrangian to be written in a single compact form with one fixed penalty
    parameter $\rho > 0$.
    \item \textbf{Single-loop stochastic primal--dual algorithm.}
    We develop a Stochastic Perturbed Proximal Augmented Lagrangian (SPPAL)
    method that combines STORM variance reduction, a damped auxiliary
    multiplier update, and a closed-form blockwise multiplier step that
    handles equality and inequality constraints in a unified way.
    \item \textbf{Residual-based convergence interface.}
    We formulate mixed-constraint KKT residuals consistent with the effective
    multiplier and establish the analysis interface needed for fixed-penalty
    complexity bounds and feasibility control in the nonlinear stochastic
    setting.
\end{itemize}

\subsection{Notation}
Throughout this paper, $\|\cdot\|$ denotes the Euclidean norm and
$\langle\cdot,\cdot\rangle$ denotes the Euclidean inner product.
For $a\in\RR$, let $[a]_+ := \max\{a,0\}$, and apply $[\cdot]_+$ to vectors
componentwise. We use $\EE_{\xi}[\cdot]$ for expectation with respect to
$\xi$, let $[p]:=\{1,\ldots,p\}$ and $\mathbb{N}:=\{0,1,2,\ldots\}$, and
write $\overline{\RR}:=\RR\cup\{+\infty\}$.
The Jacobian of $\bc$ is denoted by
$J_{\bc}(\bx):=[J_{\bc_{\calE}}(\bx),J_{\bc_{\calI}}(\bx)]\in\RR^{n\times m}$.
For $X\subseteq\RR^n$, $\delta_X$ denotes its indicator function and
$\mathrm{dist}(\bx,X):=\inf_{\by\in X}\|\bx-\by\|$.
For complexity notation, we write $K=\calO(\varepsilon)$ if
$K\leq B\varepsilon$ for some $B>0$, and $K=\Theta(\varepsilon)$ if
$b\varepsilon\leq K\leq B\varepsilon$ for some constants $0<b<B<+\infty$.

\subsection{Organization}
The rest of the paper is organized as follows. Section~\ref{sec:prelim}
introduces the assumptions and the stochastic KKT criterion.
Section~\ref{sec:method} presents the perturbed proximal ALM formulation and
the single-loop algorithm. Section~\ref{sec:analysis} establishes stochastic
complexity and residual bounds. Section~\ref{sec:experiments} reports
numerical experiments. Additional technical lemmas are provided in
Appendix~\ref{sec:appendix-supported}.
\section{Preliminaries and Assumptions}
\label{sec:prelim}

\subsection{Main Assumptions and stationarity}
\label{Subsection: main assumptions}
% \setcounter{assumption}{0}
In this subsection, we formally state our assumptions for Problem \eqref{Prob: primal problem} and the stationarity for problem.

% \begin{assumption}
%     \label{ass: proximal of nonsmooth part}
%     the function $g(\cdot)$ is proper and lower semicontinuous with effective domain denoted by  $\calX = \mathrm{dom}~g := \{x\in \R^n~|~g(x) < +\infty\}.$ Moreover, the proximal oracle of $g_i$ is available,  i.e., given $z_i \in \R^{n_i}$ and a sufficiently large constant $\mu >0$, we can solve the following problem
% 		\begin{align}\label{eq: prox-oracle-g}
% 			\min_{\bx\in \R^{n}} g_i(\bx) + \frac{\mu}{2}\|\bx-\bz\|^2.
% 		\end{align}
% \end{assumption}


\begin{assumption}
    \label{Ass: boundedness of c and f}
    Function $f(\cdot)$ is $L_f$-smooth over $\calX$.
    The mapping $\bc=(\bc_{\calE},\bc_{\calI}):\RR^n\to\RR^m$ is continuously
    differentiable over $\calX$, and there exist constants $C,G,L_c>0$ such
    that, for all $\bx,\by\in\calX$,
    \begin{equation}\label{eq: mp c const}
        \|\bc(\bx)\| \leq C,\qquad
        \|J_{\bc}(\bx)\| \leq G,\qquad
        \|J_{\bc}(\bx)-J_{\bc}(\by)\| \leq L_c\|\bx-\by\|.
    \end{equation}
\end{assumption}

\begin{assumption}
     \label{Ass: unbiased gradient oracle} 
     For any \(\bx\), there exists $\sigma > 0$ such that 
     \begin{align}
      &\EE_{\xi}[\nabla \mathbf f(\bx, \bxi)] = \nabla f(\bx), \notag \\
    & \| \nabla  \mathbf f(\bx, \bxi) - \nabla f(\bx) \| \leq \sigma, \quad \forall \xi
     \end{align}
\end{assumption}
% \begin{assumption}
%     \label{Ass: unbiased gradient oracle} 
%     There exists a $\sigma > 0$ such that 
%     \begin{align}
%     &\EE_{\xi}[\nabla \mathbf f(\bx, \bxi)] = \nabla f(\bx), \notag \\
%     & \EE_{\xi}[\| \nabla  \mathbf f(\bx, \bxi) - \nabla f(\bx) \|^{2}] \leq \sigma^{2}, \quad \forall\bx
%     \in \calX. \notag
%     \end{align}
% \end{assumption}
We also introduce a standard assumption, also known as the mean-squared smoothness assumption, widely used in variance reduction methods.
\begin{assumption}
    \label{Ass: the expected smoothness of f} 
    For almost every $\bxi \in \Xi$, $f(\cdot, \bxi)$ is differentiable and 
    for all $\bx, \by \in \calX$,
    \[
        \EE_{\bxi}\left[\left\|\nabla f(\bx; \bxi) - \nabla f(\by; \bxi)\right\|^{2}
        \right] \leq L_f^{2}\|\bx - \by\|^{2}.
    \]
\end{assumption}
\begin{assumption}[Boundedness of Multipliers]
\label{Ass: bounded multiplier}
The sequence of dual variables $\{\blmbd^k\}$ generated by 
Algorithm~\ref{alg:PLADM-eq} is contained in a convex compact set 
$\Lambda \subset \RR^{|\calE|}\times\RR_+^{|\calI|}$, i.e., there exists a
constant $B_\lambda > 0$ such that
\begin{equation}
    \|\blmbd^k\| \leq B_\lambda, \quad \forall\, k \geq 0.
\end{equation}
\end{assumption}
% Given nonconvex constraints, we impose the following nonsingularity condition on the constraints to ensure the feasibility of solutions. Constraint qualifications are often required in solving constrained optimization problems; see~\cite{shi2025momentum,lu2024variancereducedfirstordermethodsdeterministically,liu2025single,curtis2024worst}.
% \begin{assumption}
%     \label{Ass: constraints conditions} 
%     % For some sufficientlCQy small \(\varepsilon_0 > 0\), t
%     There exists constant $\delta > 0$ such that
%     \begin{align*}
%         \|\nabla \bc(\bx) \bc(\bx)\| \geq \delta \| \bc(\bx) \|, \quad \forall \bx
%     \in \RR^n.
%     \end{align*}
% \end{assumption}
\subsection{Stochastic KKT Point}
\begin{definition}[Stochastic KKT point]
\label{def:stationary}
A point $\bx^* \in \calX$ is called a \emph{stationary point} of
Problem~\eqref{Prob: primal problem} if there exists
$\blmbd^* \in \RR^{|\calE|}\times\RR_+^{|\calI|}$ such that
\begin{align}
    \bm{0}
    &\in \nabla f(\bx^*) + \partial g(\bx^*)
    + J_{\bc_{\calE}}(\bx^*)^\top \blmbd_{\calE}^*
    + J_{\bc_{\calI}}(\bx^*)^\top \blmbd_{\calI}^*
    + \mathcal{N}_{\calX}(\bx^*), \label{eq:kkt-stationarity-main}\\
    \bc_{\calE}(\bx^*) &= \bm{0}, \qquad
    \bc_{\calI}(\bx^*) \leq \bm{0}, \qquad
    \blmbd_{\calI}^* \odot \bc_{\calI}(\bx^*) = \bm{0}. \notag
\end{align}
\end{definition}

\begin{definition}[$\varepsilon$-KKT point]
\label{def:eps-kkt}
A point $\bx \in \calX$ is called an \emph{$\varepsilon$-KKT point} of
Problem~\eqref{Prob: primal problem} if there exists
$\blmbd \in \RR^{|\calE|}\times\RR_+^{|\calI|}$ such that
\begin{align}
    \mathrm{dist}\!\Bigl(
        \bm{0},\;
        \nabla f(\bx) + \partial g(\bx) + \mathcal{N}_{\calX}(\bx)
        + J_{\bc_{\calE}}(\bx)^\top \blmbd_{\calE}
        + J_{\bc_{\calI}}(\bx)^\top \blmbd_{\calI}
    \Bigr)
    &\leq \varepsilon, \notag\\
    \|\bc_{\calE}(\bx)\| \leq \varepsilon,\qquad
    \|[\bc_{\calI}(\bx)]_+\| &\leq \varepsilon, \notag\\
    \sum_{i\in\calI}|\blmbd_i c_i(\bx)| &\leq \varepsilon.
\end{align}
\end{definition}

We measure progress by the mixed-constraint residual
\begin{equation}\label{eq:kkt-residual}
    R(\bx,\blmbd)
    := \mathrm{dist}\!\Bigl(
        -\nabla f(\bx) - J_{\bc}(\bx)^\top \hat{\blmbd}(\blmbd,\bx),\;
        \partial g(\bx) + \mathcal{N}_{\calX}(\bx)
    \Bigr)^2
    + \|\bc_{\calE}(\bx)\|^2
    + \|[\bc_{\calI}(\bx)]_+\|^2.
\end{equation}
\section{Methodology and Algorithm}
\label{sec:method}
To absorb the structural difference between equality and inequality
constraints, we define the \emph{effective multiplier}
\begin{equation}\label{eq:effective-multiplier}
    \hat{\blmbd}_i(\blmbd,\bx)
    :=
    \begin{cases}
        \blmbd_i + \rho\,c_i(\bx), & i\in\calE,\\[4pt]
        \bigl[\blmbd_i + \rho\,c_i(\bx)\bigr]_+, & i\in\calI.
    \end{cases}
\end{equation}
and write $\hat{\blmbd}(\blmbd,\bx)\in\RR^m$ for the stacked vector.
This map absorbs the equality/inequality asymmetry, while the perturbation
term acts on the full multiplier vector symmetrically.

\begin{equation}\label{eq:reduced_ppl}
    \mathcal{L}_{\rho}(\bx, \blmbd, \bmu)
    =
    \underbrace{
    f(\bx)
    + \frac{1}{2\rho}\|\hat{\blmbd}(\blmbd,\bx)\|^2
    - \frac{1}{2\rho}\|\blmbd\|^2
    }_{\varphi_{\rho}(\bx,\blmbd)}
    + g(\bx)
    - \frac{1}{2\rho}\|\blmbd - \bmu\|^2.
\end{equation}

\begin{remark}
For each $i\in\calE$,
\[
    \frac{1}{2\rho}\bigl(\blmbd_i+\rho c_i(\bx)\bigr)^2
    - \frac{1}{2\rho}\blmbd_i^2
    = \blmbd_i c_i(\bx) + \frac{\rho}{2}c_i(\bx)^2.
\]
For each $i\in\calI$,
\[
    \frac{1}{2\rho}\bigl[\blmbd_i+\rho c_i(\bx)\bigr]_+^2
    - \frac{1}{2\rho}\blmbd_i^2
\]
recovers the sharp augmented-Lagrangian term for inequality constraints.
Hence $\hat{\blmbd}(\blmbd,\bx)$ absorbs the block asymmetry, while
$-\frac{1}{2\rho}\|\blmbd-\bmu\|^2$ supplies a uniform proximal perturbation.
\end{remark}

\begin{equation}\label{eq:phi-gradient}
    \nabla_{\bx}\varphi_{\rho}(\bx,\blmbd)
    = \nabla f(\bx) + J_{\bc}(\bx)^\top \hat{\blmbd}(\blmbd,\bx).
\end{equation}

\subsection{Algorithm}
\textbf{Stochastic gradient estimator.}
We replace $\nabla f(\bx^k)$ with the STORM variance-reduced estimator:
\begin{equation}\label{eq:storm}
    \bd^k =  \nabla \mathbf{f}(\bx^k, \xi^k) 
    + (1-\gamma_{k-1})\big(\bd^{k-1} - \nabla \mathbf{f}(\bx^{k-1}, \xi^k)\big),
\end{equation}where $\xi^k$ is an independent sample drawn at iteration $k$.

\textbf{Update of $\bx$.}
\begin{equation}\label{eq:x-update}
\bx^{k+1} \in \argmin_{\bx \in \calX}
\left\{
g(\bx)
+ \left\langle \hat{\nabla}_{\bx}\varphi_{\rho}(\bx^k, \blmbd^k),\ \bx - \bx^k \right\rangle
+ \frac{1}{2\eta}\|\bx - \bx^k\|^2
\right\},
\end{equation} where \(\hat{\nabla}_{\bx}\varphi_{\rho}(\bx^k, \blmbd^k)
:= \bd^k + J_{\bc}(\bx^k)^\top \hat{\blmbd}(\blmbd^k,\bx^k).\)

\textbf{Update of $\bmu$.}
\begin{equation}\label{eq:mu-update}
    \bmu^{k+1} = \bmu^k + \theta_k\big(\blmbd^k - \bmu^k\big),
\end{equation} with a diminishing step-size
\begin{equation}\label{eq:theta-def}
    \theta_k = \min\!\left\{\theta_0,\; 
    \frac{\rho\delta_k}{\|\blmbd^k - \bmu^k\|^2 + 1}\right\},
\end{equation} and \(\delta_k \) satisfies
\begin{equation}\label{eq:delta-conditions}
    \delta_0 \in (0,1], \qquad 
    \lim_{k\to\infty} \delta_k = 0, \qquad \text{and} \qquad
    \sum_{k=0}^{\infty} \delta_k = +\infty.
\end{equation}

\textbf{Update of $\blmbd$.}
\begin{equation}\label{eq:lambda-update}
    \blmbd^{k+1} = \hat{\blmbd}(\bmu^{k+1}, \bx^{k+1}).
\end{equation}
Equivalently,
\begin{align}
    \blmbd_{\calE}^{k+1}
    &= \bmu_{\calE}^{k+1} + \rho\,\bc_{\calE}(\bx^{k+1}), \notag\\
    \blmbd_{\calI}^{k+1}
    &= \bigl[\bmu_{\calI}^{k+1} + \rho\,\bc_{\calI}(\bx^{k+1})\bigr]_+.
    \notag
\end{align}
If $\blmbd_{\calI}^0,\bmu_{\calI}^0 \geq \bm{0}$, then
$\blmbd_{\calI}^k,\bmu_{\calI}^k \geq \bm{0}$ for all $k$.

\begin{algorithm}[H]
    \caption{Stochastic Perturbed Proximal Augmented Lagrangian (SPPAL)}
    \label{alg:PLADM-eq}
    \small
    \begin{algorithmic}[1]
        \STATE \textbf{Input:} $(\bx^0, \blmbd^0, \bmu^0)$ with
        $\blmbd_{\calI}^0,\bmu_{\calI}^0 \geq \bm{0}$, parameters
        $\eta>0$, $\rho>0$, $\theta_0\in(0,1]$, momentum sequence
        $\{\gamma_k\}$, and damping sequence $\{\delta_k\}$ satisfying
        $\delta_k\to 0$ and $\sum_k \delta_k = +\infty$.
        \FOR{$k = 0, 1, 2, \ldots$}
            \STATE Update $\bx^{k+1}$ by \eqref{eq:x-update}.
            \STATE Update $\bmu^{k+1}$ by \eqref{eq:mu-update}.
            \STATE Update $\blmbd^{k+1}$ by \eqref{eq:lambda-update}.
        \ENDFOR
        \STATE \textbf{Output:} $\bx^R$ (randomly selected iterate).
    \end{algorithmic}
\end{algorithm}
\section{Stochastic Complexity Analysis}
\label{sec:analysis}
\begin{lemma}[Smoothness of $\varphi_{\rho}$]
\label{lem:phi-smooth}
Under Assumptions~\ref{Ass: boundedness of c and f} 
and~\ref{Ass: bounded multiplier}, the function 
$\varphi_{\rho}(\cdot,\blmbd)$ is $L_{\varphi}$-smooth for any 
$\blmbd$ with $\|\blmbd\|\leq B_{\lambda}$, i.e.,
\begin{equation}\label{eq:phi-smooth}
    \|\nabla_{\bx}\varphi_{\rho}(\bx,\blmbd)
     - \nabla_{\bx}\varphi_{\rho}(\by,\blmbd)\|
    \leq L_{\varphi}\|\bx-\by\|,
    \quad \forall\,\bx,\by\in\calX,
\end{equation}
where
\begin{equation}\label{eq:Lphi-def}
    L_{\varphi} := L_f + L_c(B_{\lambda}+\rho C) + \rho G^2.
\end{equation}
\end{lemma}

% \begin{proof}
% Recall that
% \(
% \nabla_{\bx}\varphi_{\rho}(\bx,\blmbd)
% = \nabla f(\bx) + J_{\bc}(\bx)^{\top}\!\bigl(\blmbd+\rho\,\bc(\bx)\bigr).
% \)
% For any $\bx,\by\in\RR^n$, add and subtract 
% $J_{\bc}(\bx)^{\top}(\blmbd+\rho\,\bc(\by))$ to obtain
% \begin{align}
% &\|\nabla_{\bx}\varphi_{\rho}(\bx,\blmbd)
%   -\nabla_{\bx}\varphi_{\rho}(\by,\blmbd)\| \notag\\
% &\leq
%   \underbrace{\|\nabla f(\bx)-\nabla f(\by)\|}_{\mathrm{(I)}}
%   +\underbrace{\bigl\|[J_{\bc}(\bx)-J_{\bc}(\by)]^{\top}
%     (\blmbd+\rho\,\bc(\by))\bigr\|}_{\mathrm{(II)}}
%   +\underbrace{\bigl\|J_{\bc}(\bx)^{\top}\rho
%     [\bc(\bx)-\bc(\by)]\bigr\|}_{\mathrm{(III)}}.
%     \label{eq:phi-three-terms}
% \end{align}

% \medskip\noindent
% \textbf{Term (I).}
% By $L_f$-smoothness of $f$:
% \begin{equation}
%     \mathrm{(I)} \leq L_f\|\bx-\by\|.
% \end{equation}

% \medskip\noindent
% \textbf{Term (II).}
% Since $\bc$ is $L_c$-smooth, we have 
% $\|J_{\bc}(\bx)-J_{\bc}(\by)\|\leq L_c\|\bx-\by\|$.
% Using $\|\blmbd\|\leq B_{\lambda}$ and $\|\bc(\by)\|\leq C$ 
% from Assumptions~\ref{Ass: bounded multiplier} 
% and~\ref{Ass: boundedness of c and f}:
% \begin{equation}
%     \mathrm{(II)}
%     \leq L_c\|\bx-\by\|\cdot\|\blmbd+\rho\,\bc(\by)\|
%     \leq L_c(B_{\lambda}+\rho C)\|\bx-\by\|.
% \end{equation}

% \medskip\noindent
% \textbf{Term (III).}
% By $\|J_{\bc}(\bx)\|\leq G$ and the $G$-Lipschitz 
% continuity of $\bc$ (Assumption~\ref{Ass: boundedness of c and f}):
% \begin{equation}
%     \mathrm{(III)}
%     \leq \rho\|J_{\bc}(\bx)\|\cdot\|\bc(\bx)-\bc(\by)\|
%     \leq \rho G^2\|\bx-\by\|.
% \end{equation}

% \medskip\noindent
% Substituting (I)--(III) back into \eqref{eq:phi-three-terms} gives
% \[
%     \|\nabla_{\bx}\varphi_{\rho}(\bx,\blmbd)
%      -\nabla_{\bx}\varphi_{\rho}(\by,\blmbd)\|
%     \leq \bigl(L_f + L_c(B_{\lambda}+\rho C)+\rho G^2\bigr)
%          \|\bx-\by\|
%     = L_{\varphi}\|\bx-\by\|,
% \]
% which completes the proof.
% \end{proof}
\begin{lemma}[Basic relations on $\blmbd$ and $\bmu$]
\label{lem:basic-relations}
Let $\{(\bx^k, \blmbd^k, \bmu^k)\}$ be the sequence generated by 
Algorithm~\ref{alg:PLADM-eq}. Under Assumption~\ref{Ass: boundedness of c and f}, the following 
relations hold for all $k \geq 0$:
\begin{align}
    \|\bmu^{k+1} - \bmu^k\|^2 = \theta_k^2 \|\blmbd^{k} - \bmu^k\|^2  &\leq \frac{\rho^2\delta_k^2}{4}, 
    \label{eq:mu-diff}\\
    \theta_k \|\blmbd^k - \bmu^k\|^2 &\leq \rho\delta_k, 
    \label{eq:theta-bound}\\
    \|\bmu^{k+1} - \blmbd^k\|^2 &= (1-\theta_k)^2\|\blmbd^k - \bmu^k\|^2,
    \label{eq:mu-lambda-diff}\\
    \mathbb{E}[\|\blmbd^{k+1} - \blmbd^k\|^2] &\leq 
    2\rho^2 L_c^2\,\mathbb{E}[\|\bx^{k+1} - \bx^k\|^2] + \frac{\rho^2\delta_k^2}{2}.
    \label{eq:lambda-diff}
\end{align}
\end{lemma}
% \begin{proof}
% \textbf{Relation \eqref{eq:mu-diff}.}
% From the $\bmu$-update \eqref{eq:mu-update},
% $\bmu^{k+1} - \bmu^k = \theta_k(\blmbd^k - \bmu^k)$, so
% \begin{align*}
%   \|\bmu^{k+1} - \bmu^k\|^2
%   = \theta_k^2\|\blmbd^k - \bmu^k\|^2
%   \leq \frac{\rho^2\delta_k^2}{\bigl(\|\blmbd^k-\bmu^k\|^2+1\bigr)^2}
%     \cdot \|\blmbd^k-\bmu^k\|^2
%   \leq \frac{\rho^2\delta_k^2}{4},
% \end{align*}
% where the last inequality uses $\frac{t}{(t+1)^2} \leq \frac{1}{4}$ for $t = \|\blmbd^k-\bmu^k\|^2 \geq 0$.
 
% \textbf{Relation \eqref{eq:theta-bound}.}
% By definition \eqref{eq:theta-def},
% $\theta_k \leq \frac{\rho\delta_k}{\|\blmbd^k-\bmu^k\|^2+1}$, hence
% \begin{align*}
%   \theta_k\|\blmbd^k-\bmu^k\|^2
%   \leq \frac{\rho\delta_k\|\blmbd^k-\bmu^k\|^2}{\|\blmbd^k-\bmu^k\|^2+1}
%   \leq \rho\delta_k.
% \end{align*}
 
% \textbf{Relation \eqref{eq:mu-lambda-diff}.}
% $\bmu^{k+1} - \blmbd^k = \bmu^k - \blmbd^k + \theta_k(\blmbd^k-\bmu^k)
% = (1-\theta_k)(\bmu^k - \blmbd^k)$,
% so $\|\bmu^{k+1}-\blmbd^k\|^2 = (1-\theta_k)^2\|\blmbd^k-\bmu^k\|^2$.
 
% \textbf{Relation \eqref{eq:lambda-diff}.}
% From the $\blmbd$-update \eqref{eq:lambda-update},
% $\blmbd^{k+1} - \blmbd^k
% = (\bmu^{k+1}-\bmu^k) + \rho\bigl(\bc(\bx^{k+1})-\bc(\bx^k)\bigr)$.
% Using $(a+b)^2 \leq 2(a^2+b^2)$ and $L_c$-smoothness of $\bc$:
% \begin{align*}
%   \|\blmbd^{k+1}-\blmbd^k\|^2
%   &\leq 2\|\bmu^{k+1}-\bmu^k\|^2 + 2\rho^2\|\bc(\bx^{k+1})-\bc(\bx^k)\|^2 \\
%   &\leq \frac{\rho^2\delta_k^2}{2} + 2\rho^2 L_c^2\|\bx^{k+1}-\bx^k\|^2.
% \end{align*}
% Taking expectation gives \eqref{eq:lambda-diff}.
% \end{proof}
% \begin{lemma}[Boundedness of multipliers]
% \label{lem:bounded-multiplier}
% Let $\{(\bx^k, \blmbd^k, \bmu^k)\}$ be the sequence generated by 
% Algorithm~\ref{alg:PLADM-eq}. 
% Under Assumption~\ref{Ass: boundedness of c and f}, for all $k \geq 0$,  it holds that
% \begin{align}
%     \|\bmu^{k+1}\| &\leq \|\bmu^0\| + \rho C\sum_{j=0}^{k}\tau_j\theta_j, 
%     \label{eq:mu-bound}\\
%     \|\blmbd^{k+1}\| &\leq \|\bmu^0\| + \rho C\left(\tau_k + \sum_{j=0}^{k}\tau_j\theta_j\right).
%     \label{eq:lambda-bound}
% \end{align}
% \end{lemma}

% \begin{lemma}[Primal Feasibility bound]
% \label{lem:bounded-multiplier}
% Let $\{(\bx^k, \blmbd^k, \bmu^k)\}$ be the sequence generated by 
% Algorithm~\ref{alg:PLADM-eq} with \(\tau_k = \tau\). Under Assumption~\ref{Ass: boundedness of c and f}, for all $k \geq 0$,  it holds that
% \begin{align}
%  \frac{1}{T}\sum_{k = 0}^{T-1}\mathbb{E}\big[\|\bc(\bx^{k+1})\|^2\big]  = \mathcal{O}(logT/T)\big]
% \end{align} \blue{with \(\delta_k = \frac{c}{k+1}\)}.

% \end{lemma}
% \begin{proof}
%     \begin{align}
%  \frac{1}{T}\sum_{k = 0}^{T-1}\mathbb{E}\big[\|\bc(\bx^{k+1})\|^2\big]  = \frac{1}{\rho^2\tau^2}\frac{1}{T}\sum_{k = 0}^{T-1}\mathbb{E}\big[\|\blmbd^{k+1}-\bmu^{k+1}\|^2\big].
% \end{align}
% \end{proof}

% \begin{lemma}[sufficient descent]
% \end{lemma}


% \begin{lemma}[primal feasibility bound]
% \end{lemma}

% ================================================================
%  Sufficient Descent + Primal Feasibility  (equality-constrained)
%  Adapted from PLADA Lemma 33 to stochastic / no-slack setting
% ================================================================

% ----------------------------------------------------------------
% Notation reminders (consistent with contrained_format__1_.tex):
%   L_rho(x,lambda,mu) = f(x)+g(x)+<lambda,c(x)>
%                        + rho/2*||c(x)||^2 - 1/(2rho)*||lambda-mu||^2
%   x-update  : proximal-gradient with STORM estimator d^k
%   mu-update : mu^{k+1} = mu^k + theta_k(lambda^k - mu^k),
%               theta_k = delta_k / (1 + ||lambda^k - mu^k||^2)
%   lambda-update (tau_k = tau, constant):
%               lambda^{k+1} = mu^{k+1} + rho*tau*c(x^{k+1})
%
%  Key consequence of lambda-update (used everywhere):
%   lambda^{k+1} - mu^{k+1} = rho*tau*c(x^{k+1})           ... (*)
% ----------------------------------------------------------------
\begin{lemma}[Sufficient Descent]
\label{lem:sufficient-descent}
Let $\{(\bx^k, \blmbd^k, \bmu^k)\}$ be the sequence generated by
Algorithm~\ref{alg:PLADM-eq}. Under
Assumptions~\ref{Ass: boundedness of c and f},
\ref{Ass: unbiased gradient oracle}, and~\ref{Ass: bounded multiplier}, the following inequality holds for all $k \ge 0$:
\begin{align}
  \mathbb{E}\!\left[\calL_{\rho}(\bx^{k+1}, \blmbd^{k+1}, \bmu^{k+1})
  \,\big|\, \mathcal{F}_k\right]
  &\;\le\;
  \calL_{\rho}(\bx^k, \blmbd^k, \bmu^k)
  - C_1\,\mathbb{E}\!\left[\|\bx^{k+1}-\bx^k\|^2
    \,\big|\, \mathcal{F}_k\right]
  + \frac{\eta}{2}\|\be^k\|^2
  + \hat\delta_k,
  \label{eq:sufficient-descent}
\end{align}
where
\begin{equation}\label{eq:C1-deltahat-def}
  C_1 := \frac{1}{2\eta} - \frac{L_\varphi}{2} - \rho L_c^2 > 0,
  \qquad
  \hat\delta_k := \frac{\rho\delta_k^2}{4} + \delta_k.
\end{equation}
Furthermore, rearranging \eqref{eq:sufficient-descent} gives
\begin{align}\label{eq:descent-rearranged}
  C_1\,\mathbb{E}\!\left[\|\bx^{k+1}-\bx^k\|^2 \mid \mathcal{F}_k\right]
  \leq\;
  &\calL_{\rho}(\bx^k, \blmbd^k, \bmu^k)
  - \mathbb{E}\!\left[\calL_{\rho}(\bx^{k+1}, \blmbd^{k+1}, \bmu^{k+1})
    \mid \mathcal{F}_k\right]
  + \frac{\eta}{2}\|\be^k\|^2
  + \hat\delta_k.
\end{align}
\end{lemma}

% \begin{proof}
% Decompose the one-step change of $\calL_\rho$ as
% \begin{align}
%   &\calL_\rho(\bx^{k+1},\blmbd^{k+1},\bmu^{k+1})
%    - \calL_\rho(\bx^k,\blmbd^k,\bmu^k) \notag\\
%   &= \underbrace{
%       \calL_\rho(\bx^{k+1},\blmbd^k,\bmu^k)
%     - \calL_\rho(\bx^k,\blmbd^k,\bmu^k)
%     }_{\text{Part A}}
%   + \underbrace{
%       \calL_\rho(\bx^{k+1},\blmbd^{k+1},\bmu^{k+1})
%     - \calL_\rho(\bx^{k+1},\blmbd^k,\bmu^k)
%     }_{\text{Part B}}.
%   \label{eq:decompose}
% \end{align}

% \medskip\noindent\textbf{Part A (x-update).}
% By the $L_\varphi$-smoothness of $\varphi_\rho(\cdot,\blmbd^k)$
% (Lemma~\ref{lem:phi-smooth}) and the descent lemma:
% \begin{align}
%   \text{Part A}
%   &\leq \langle \nabla_\bx\varphi_\rho(\bx^k,\blmbd^k),\, \bx^{k+1}-\bx^k\rangle
%        + g(\bx^{k+1}) - g(\bx^k)
%        + \frac{L_\varphi}{2}\|\bx^{k+1}-\bx^k\|^2.
%   \label{eq:partA-descent}
% \end{align}
% From the $\bx$-update optimality condition
% $-\hat{\nabla}_\bx\varphi_\rho(\bx^k,\blmbd^k)
%  - \frac{1}{\eta}(\bx^{k+1}-\bx^k) \in \partial g(\bx^{k+1})$,
% taking $\bx = \bx^k$ in the subgradient inequality gives
% \begin{equation}
%   g(\bx^{k+1}) - g(\bx^k)
%   \leq \Bigl\langle
%     -\hat{\nabla}_\bx\varphi_\rho(\bx^k,\blmbd^k)
%     - \tfrac{1}{\eta}(\bx^{k+1}-\bx^k),\;
%     \bx^{k+1}-\bx^k
%   \Bigr\rangle.
%   \label{eq:partA-sub}
% \end{equation}
% Substituting \eqref{eq:partA-sub} into \eqref{eq:partA-descent}
% and using $\hat{\nabla}_\bx\varphi_\rho(\bx^k,\blmbd^k)
% = \nabla_\bx\varphi_\rho(\bx^k,\blmbd^k) - \be^k$:
% \begin{align}
%   \text{Part A}
%   &\leq \langle \be^k,\, \bx^{k+1}-\bx^k\rangle
%        + \Bigl(\frac{L_\varphi}{2} - \frac{1}{2\eta}\Bigr)
%          \|\bx^{k+1}-\bx^k\|^2.
% \end{align}
% Taking conditional expectation $\mathbb{E}[\,\cdot\,|\mathcal{F}_k]$
% and applying Young's inequality
% $\langle \be^k, \bx^{k+1}-\bx^k\rangle
%  \leq \frac{\eta}{2}\|\be^k\|^2
%     + \frac{1}{2\eta}\|\bx^{k+1}-\bx^k\|^2$:
% \begin{align}
%   \mathbb{E}[\text{Part A}\,|\,\mathcal{F}_k]
%   &\leq -\Bigl(\frac{1}{2\eta} - \frac{L_\varphi}{2}\Bigr)
%         \mathbb{E}[\|\bx^{k+1}-\bx^k\|^2\,|\,\mathcal{F}_k]
%        + \frac{\eta}{2}\|\be^k\|^2.
%   \label{eq:partA-bound}
% \end{align}

% \medskip\noindent\textbf{Part B (dual-variable change).}
% Expanding $\calL_\rho$ using definition \eqref{eq:reduced_ppl}:
% \begin{align}
%   \text{Part B}
%   &= \langle \blmbd^{k+1}-\blmbd^k,\, \bc(\bx^{k+1})\rangle
%     - \frac{1}{2\rho}\Bigl(
%         \|\blmbd^{k+1}-\bmu^{k+1}\|^2
%       - \|\blmbd^k-\bmu^k\|^2
%     \Bigr).
%   \label{eq:partB-expand}
% \end{align}
% Substituting
% $\bc(\bx^{k+1}) = \frac{1}{\rho}(\blmbd^{k+1}-\bmu^{k+1})$
% from \eqref{eq:dual-gap-feasibility} into the inner product, and applying
% the polarization identity
% $\langle a,b\rangle = \frac{1}{2}\|a\|^2+\frac{1}{2}\|b\|^2
% -\frac{1}{2}\|a-b\|^2$
% with $a = \blmbd^{k+1}-\blmbd^k$ and $b = \blmbd^{k+1}-\bmu^{k+1}$
% (so $a - b = \bmu^{k+1}-\blmbd^k$):
% \begin{align}
%   \text{Part B}
%   &= \frac{1}{2\rho}\|\blmbd^{k+1}-\blmbd^k\|^2
%     - \frac{1}{2\rho}\|\bmu^{k+1}-\blmbd^k\|^2
%     + \frac{1}{2\rho}\|\blmbd^k-\bmu^k\|^2,
%   \label{eq:partB-identity}
% \end{align}
% where $\|\blmbd^{k+1}-\bmu^{k+1}\|^2$ cancels exactly.
% We now bound each term using Lemma~\ref{lem:basic-relations}.

% \medskip\noindent
% \textit{Second and third terms.}
% By \eqref{eq:mu-lambda-diff},
% $\|\bmu^{k+1}-\blmbd^k\|^2 = (1-\theta_k)^2\|\blmbd^k-\bmu^k\|^2$,
% so
% \begin{equation}
%   \frac{1}{2\rho}\|\blmbd^k-\bmu^k\|^2
%   - \frac{1}{2\rho}\|\bmu^{k+1}-\blmbd^k\|^2
%   = \frac{2\theta_k-\theta_k^2}{2\rho}\|\blmbd^k-\bmu^k\|^2
%   \leq \frac{\theta_k}{\rho}\|\blmbd^k-\bmu^k\|^2.
% \end{equation}
% By \eqref{eq:theta-bound}, $\theta_k\|\blmbd^k-\bmu^k\|^2 \leq \rho\delta_k$,
% hence this is $\leq \delta_k$.

% \medskip\noindent
% \textit{First term.}
% Taking conditional expectation and applying \eqref{eq:lambda-diff}:
% \begin{equation}
%   \frac{1}{2\rho}
%   \mathbb{E}[\|\blmbd^{k+1}-\blmbd^k\|^2\,|\,\mathcal{F}_k]
%   \leq \rho L_c^2\,
%   \mathbb{E}[\|\bx^{k+1}-\bx^k\|^2\,|\,\mathcal{F}_k]
%   + \frac{\rho\delta_k^2}{4}.
% \end{equation}

% \medskip\noindent
% Combining the two estimates:
% \begin{align}
%   \mathbb{E}[\text{Part B}\,|\,\mathcal{F}_k]
%   \leq \rho L_c^2\,
%        \mathbb{E}[\|\bx^{k+1}-\bx^k\|^2\,|\,\mathcal{F}_k]
%        + \frac{\rho\delta_k^2}{4} + \delta_k.
%   \label{eq:partB-bound}
% \end{align}

% \medskip\noindent\textbf{Combining Parts A and B.}
% Adding \eqref{eq:partA-bound} and \eqref{eq:partB-bound}:
% \begin{align}
%   &\mathbb{E}[\calL_\rho(\bx^{k+1},\blmbd^{k+1},\bmu^{k+1})\,|\,\mathcal{F}_k]
%    - \calL_\rho(\bx^k,\blmbd^k,\bmu^k) \notag\\
%   &\leq -\underbrace{\Bigl(
%       \frac{1}{2\eta} - \frac{L_\varphi}{2} - \rho L_c^2
%     \Bigr)}_{=\,C_1}
%     \mathbb{E}[\|\bx^{k+1}-\bx^k\|^2\,|\,\mathcal{F}_k]
%    + \frac{\eta}{2}\|\be^k\|^2
%    + \underbrace{\frac{\rho\delta_k^2}{4} + \delta_k}_{=\,\hat\delta_k}.
% \end{align}
% Condition \eqref{eq:eta-condition} ensures $C_1 > 0$,
% which completes the proof.
% \end{proof}

\begin{lemma}[Primal Stationarity Residual Bound]
\label{lem:stationarity}
Let $\{(\bx^k, \blmbd^k, \bmu^k)\}$ be the sequence generated by
Algorithm~\ref{alg:PLADM-eq}.
Under Assumptions~\ref{Ass: boundedness of c and f}
and~\ref{Ass: unbiased gradient oracle},
define the stationarity residual
\begin{equation}
    \bzeta_{\bx}^{k+1} \in \nabla_{\bx}\varphi_\rho(\bx^{k+1}, \blmbd^{k+1})
    + \partial g(\bx^{k+1}) + \mathcal{N}_{\calX}(\bx^{k+1}).
\end{equation}
Then for all $k \geq 0$,
% \begin{equation}\label{eq:zeta-bound}
%     \|\bzeta_{\bx}^{k+1}\|^2
%     \;\leq\;
%     4\!\left(L_\varphi + \frac{1}{\eta}\right)^{\!2}
%     \|\bx^{k+1} - \bx^k\|^2
%     \;+\; 4\|\be^k\|^2
%     \;+\; 4G^2\|\blmbd^{k+1} - \blmbd^k\|^2,
% \end{equation}
% Consequently, taking expectation and applying
Lemma~\ref{lem:basic-relations}~\eqref{eq:lambda-diff},
\begin{equation}\label{eq:zeta-expect}
    \mathbb{E}\!\left[\|\bzeta_{\bx}^{k+1}\|^2\right]
    \;\leq\;
    4\!\left[\left(L_\varphi+\frac{1}{\eta}\right)^{\!2}
             + 2G^2\rho^2 L_c^2\right]
    \mathbb{E}\!\left[\|\bx^{k+1}-\bx^k\|^2\right]
    \;+\; 4\,\mathbb{E}\!\left[\|\be^k\|^2\right]
    \;+\; 2G^2\rho^2\delta_k^2.
\end{equation} where $\be^k := \bd^k - \nabla f(\bx^k)$ is the STORM estimation error.
\end{lemma}


 
% \begin{proof}
 
% \medskip\noindent\textbf{Step 1: Residual decomposition.}
% From the first-order optimality condition of the $\bx$-update
% \eqref{eq:x-update}, there exists $v^{k+1} \in \partial g(\bx^{k+1})$
% such that
% \begin{equation}\label{eq:x-opt}
%     v^{k+1}
%     = -\bd^k
%       - J_{\bc}(\bx^k)^{\top}\!\left(\blmbd^k + \rho\,\bc(\bx^k)\right)
%       - \frac{1}{\eta}\left(\bx^{k+1} - \bx^k\right).
% \end{equation}
% Setting $\bzeta^{k+1}
% := \nabla f(\bx^{k+1}) + v^{k+1}
%    + J_{\bc}(\bx^{k+1})^{\top}\blmbd^{k+1}$
% and substituting \eqref{eq:x-opt},
% we add and subtract $\nabla_{\bx}\varphi_{\rho}(\bx^k, \blmbd^{k+1})$
% to obtain
% \begin{align}
%     \bzeta^{k+1}
%     &= \underbrace{
%          \nabla_{\bx}\varphi_{\rho}(\bx^{k+1}, \blmbd^{k+1})
%          - \nabla_{\bx}\varphi_{\rho}(\bx^k, \blmbd^{k+1})
%        }_{\displaystyle\mathrm{(I)}}
%     + \underbrace{
%          \nabla_{\bx}\varphi_{\rho}(\bx^k, \blmbd^{k+1})
%          - \nabla_{\bx}\varphi_{\rho}(\bx^k, \blmbd^k)
%        }_{\displaystyle\mathrm{(II)}}
%     - \be^k
%     - \frac{1}{\eta}\left(\bx^{k+1} - \bx^k\right).
%     \label{eq:zeta-decomp}
% \end{align}
% Note that by definition of $\varphi_\rho$,
% \begin{equation}\label{eq:phi-lambda-diff}
%     \nabla_{\bx}\varphi_{\rho}(\bx, \blmbd_1)
%     - \nabla_{\bx}\varphi_{\rho}(\bx, \blmbd_2)
%     = J_{\bc}(\bx)^{\top}(\blmbd_1 - \blmbd_2),
% \end{equation}
% so $\mathrm{(II)} = J_{\bc}(\bx^k)^{\top}(\blmbd^{k+1} - \blmbd^k)$.
 
% \medskip\noindent\textbf{Step 2: Bounding each term.}
% Applying $(a_1+a_2+a_3+a_4)^2 \leq 4(a_1^2+a_2^2+a_3^2+a_4^2)$
% to \eqref{eq:zeta-decomp}:
 
% \medskip\noindent\textit{Term~(I).}
% By $L_\varphi$-smoothness of $\varphi_\rho(\cdot,\blmbd^{k+1})$
% (Lemma~\ref{lem:phi-smooth}):
% \begin{equation}
%     \|\mathrm{(I)}\|
%     \;\leq\; L_\varphi\|\bx^{k+1} - \bx^k\|.
% \end{equation}
 
% \medskip\noindent\textit{Term~(II).}
% By \eqref{eq:phi-lambda-diff} and
% Assumption~\ref{Ass: boundedness of c and f}:
% \begin{equation}
%     \|\mathrm{(II)}\|
%     \;=\; \|J_{\bc}(\bx^k)^{\top}(\blmbd^{k+1} - \blmbd^k)\|
%     \;\leq\; G\|\blmbd^{k+1} - \blmbd^k\|.
% \end{equation}
 
% \medskip\noindent
% Collecting the four terms and squaring yields \eqref{eq:zeta-bound}.
 
% \medskip\noindent\textbf{Step 3: Taking expectation.}
% Applying Lemma~\ref{lem:basic-relations}~\eqref{eq:lambda-diff}:
% \begin{equation}
%     4G^2\,\mathbb{E}\!\left[\|\blmbd^{k+1}-\blmbd^k\|^2\right]
%     \;\leq\;
%     8G^2\rho^2 L_c^2\,\mathbb{E}\!\left[\|\bx^{k+1}-\bx^k\|^2\right]
%     + 2G^2\rho^2\delta_k^2.
% \end{equation}
% Substituting into \eqref{eq:zeta-bound} and taking expectation
% yields \eqref{eq:zeta-expect}.
% \end{proof}

% \begin{lemma}[Primal Feasibility Bound]
% \label{lem:primal-feasibility}
% Under the same conditions as Lemma~\ref{lem:sufficient-descent},
% with $\delta_k = c/(k+1)$ for some $c > 0$, we have
% \begin{align}
%   \frac{1}{T}\sum_{k=0}^{T-1}\mathbb{E}\!\left[\|\bc(\bx^{k+1})\|^2\right]
%   \;=\; \mathcal{O}\!\left(\frac{\log T}{T}\right).
%   \label{eq:feasibility-rate}
% \end{align}
% \end{lemma}
% \begin{proof}
% \textbf{Step 1: Derive the residual expression.}
% From the first-order optimality condition of the $\bx$-update 
% \eqref{eq:x-update}, there exists $v^{k+1} \in \partial g(\bx^{k+1})$ 
% such that
% \begin{equation}\label{eq:x-opt}
%     v^{k+1} = -\bd^k - J_{\bc}(\bx^k)^\top\!\left(\blmbd^k 
%     + \rho\,\bc(\bx^k)\right) - \frac{1}{\eta}(\bx^{k+1} - \bx^k).
% \end{equation}
% Setting $\bzeta^{k+1} := \nabla f(\bx^{k+1}) + v^{k+1} 
% + J_{\bc}(\bx^{k+1})^\top\blmbd^{k+1}$ and substituting 
% \eqref{eq:x-opt}, we add and subtract 
% $\nabla_{\bx}\varphi_{\rho}(\bx^{k+1}, \blmbd^k)$ to obtain
% \begin{align}
%     \bzeta^{k+1} 
%     &= \underbrace{\bigl[\nabla_{\bx}\varphi_{\rho}(\bx^{k+1},\blmbd^k)
%       - \nabla_{\bx}\varphi_{\rho}(\bx^k,\blmbd^k)\bigr]}_{\text{(I)}}
%     - \be^k
%     + \underbrace{J_{\bc}(\bx^{k+1})^\top
%       (\bmu^{k+1} - \blmbd^k)}_{\text{(II)}}
%     - \frac{1}{\eta}(\bx^{k+1} - \bx^k),
%     \label{eq:zeta-decomp}
% \end{align}
% where we used the $\blmbd$-update \eqref{eq:lambda-update} to simplify:
% \begin{equation}
%     \blmbd^{k+1} - \blmbd^k - \rho\,\bc(\bx^{k+1}) 
%     = \bmu^{k+1} + \rho\,\bc(\bx^{k+1}) - \blmbd^k 
%       - \rho\,\bc(\bx^{k+1}) 
%     = \bmu^{k+1} - \blmbd^k.
% \end{equation}

% \textbf{Step 2: Bound each term.}
% Applying $(a+b+c+d)^2 \leq 4(a^2+b^2+c^2+d^2)$ to 
% \eqref{eq:zeta-decomp}:

% \medskip\noindent\textit{Term (I).} 
% By $L_\varphi$-smoothness of $\varphi_\rho(\cdot,\blmbd^k)$ with 
% $\blmbd^k$ fixed:
% \begin{equation}
%     \|\text{(I)}\| \leq L_\varphi\|\bx^{k+1} - \bx^k\|.
% \end{equation}

% \medskip\noindent\textit{Term (II).} 
% By Assumption~\ref{Ass: boundedness of c and f} and 
% Lemma~\ref{lem:basic-relations} \eqref{eq:mu-lambda-diff}:
% \begin{equation}
%     \|\text{(II)}\| \leq G\|\bmu^{k+1} - \blmbd^k\| 
%     = G(1-\theta_k)\|\blmbd^k - \bmu^k\|.
% \end{equation}

% \medskip\noindent
% Combining the four terms and applying $(a+b+c+d)^2\leq 4(a^2+b^2+c^2+d^2)$
% yields \eqref{eq:zeta-bound} with constant absorbed into $C_2$.

% \textbf{Step 3: Ergodic bound.}
% Taking expectation and averaging over $k = 0,\ldots,T-1$:
% \begin{align}
%     \frac{1}{T}\sum_{k=0}^{T-1}\mathbb{E}[\|\bzeta^{k+1}\|^2]
%     &\leq \frac{C_2}{T}\sum_{k=0}^{T-1}
%       \mathbb{E}[\|\bx^{k+1}-\bx^k\|^2]
%     + \frac{1}{T}\sum_{k=0}^{T-1}\mathbb{E}[\|\be^k\|^2]
%     + \frac{3G^2}{T}\sum_{k=0}^{T-1}
%       \mathbb{E}[\|\bmu^{k+1}-\blmbd^k\|^2].
% \end{align}
% % The three terms are bounded as follows:
% % \begin{itemize}
% %     \item \textbf{First term:} By Lemma~\ref{lem:sufficient-descent} 
% %     \eqref{eq:sufficient-descent} and telescoping, 
% %     $\frac{1}{T}\sum\mathbb{E}[\|\bx^{k+1}-\bx^k\|^2] 
% %     = \mathcal{O}(\log T/T)$.
% %     \item \textbf{Second term:} By the STORM variance bound 
% %     (Lemma~\ref{lem:storm}), 
% %     $\frac{1}{T}\sum\mathbb{E}[\|\be^k\|^2] 
% %     = \mathcal{O}(\log T / T)$.
% %     \item \textbf{Third term:} By Lemma~\ref{lem:basic-relations} 
% %     \eqref{eq:mu-lambda-diff} and \eqref{eq:theta-bound},
% %     $\|\bmu^{k+1}-\blmbd^k\|^2 = (1-\theta_k)^2\|\blmbd^k-\bmu^k\|^2$,
% %     and $\theta_k\|\blmbd^k-\bmu^k\|^2 \leq \rho\delta_k$, so
% %     $\frac{1}{T}\sum\mathbb{E}[\|\bmu^{k+1}-\blmbd^k\|^2] 
% %     = \mathcal{O}(\log T/T)$.
% % \end{itemize}
% Combining the three bounds gives \eqref{eq:ergodic-stationarity}.
% \end{proof}

\begin{lemma}[Recursive Gradient Error]
\label{lem:storm}
Suppose that Assumptions~\ref{Ass: boundedness of c and f},
\ref{Ass: unbiased gradient oracle}, and
\ref{Ass: the expected smoothness of f} hold.
Then, for any $k \geq 0$, it holds that
\begin{equation}\label{eq:storm-error-bound}
    \EE\!\left[\|\be^{k+1}\|^2\right]
    \leq (1 - \gamma_k)^2\,\EE\!\left[\|\be^k\|^2\right]
    + 2(1 - \gamma_k)^2L_f^2\,\EE\!\left[\|\bx^{k+1} - \bx^k\|^2\right]
    + 2\gamma_k^2\sigma^2.
\end{equation}
% where $\be^k := \bd^k - \nabla f(\bx^k)$ is the STORM estimation error.
\end{lemma}

\begin{lemma}\label{lm: the variance reduction}
Under Assumptions~\ref{Ass: boundedness of c and f}--\ref{Ass: the expected
smoothness of f}, let $\{(\bx^k,\blmbd^k,\bmu^k)\}$ be generated by
Algorithm~1 with non-increasing step sizes $\{\eta_k\}$ satisfying
% $\eta_1 := \eta_0 = \sup_k \eta_k$ and
\begin{equation}\label{eq:step-cond}
    % \eta_k < \frac{1}{2(L_\varphi + 2\rho L_c^2)},
    % \qquad
    % \eta_k \leq \frac{\sqrt{\gamma_k(2-\gamma_k)}}{2L_f},
    % \quad \forall\, k \geq 0.
    (1 - L_{c,\varphi} \eta_k)^{-1} \leq 2
\end{equation} where \(L_{c,\varphi} = L_{\varphi} + 2 \rho L_c^2\). Then for any $K \geq 1$,
\begin{equation}\label{eq:main-bound}
    \frac{1}{K}\sum_{k=0}^{K-1}\gamma_k\,\EE\!\left[\|\be^k\|^2\right]
    \leq
    \frac{16\eta_1 L_f^2\bigl(\calL_\rho^0 - \calL_\rho^*\bigr)}{K}
    + \frac{2\,\EE\!\left[\|\be^1\|^2\right]}{K}
    + \frac{16\eta_1 L_f^2}{K}\sum_{k=0}^{K-1}\hat{\delta}_k
    + \frac{1}{K}\sum_{k=0}^{K-1}4\gamma_k^2\sigma^2,
\end{equation}
where $\calL_\rho^* \geq F^*$ is a lower bound on $\calL_\rho$.
\end{lemma}

% \begin{proof}

% \medskip\noindent\textbf{Step 1: Lyapunov function.}
% Define
% \begin{equation}\label{eq:Phi-def}
%     \varphi^k \;:=\; \EE\!\left[\calL_\rho(\bx^k,\blmbd^k,\bmu^k)\right]
%     + \frac{1}{8\eta_1 L_f^2}\,\EE\!\left[\|\be^k\|^2\right].
% \end{equation}
% Since $\eta_1 \geq \eta_k$ for all $k$, the coefficient $\frac{1}{8\eta_1 L_f^2}$
% is fixed.

% \medskip\noindent\textbf{Step 2: One-step descent.}
% From Lemma~\ref{lem:descent} (taking full expectation of~\eqref{eq:descent}):
% \begin{equation}\label{eq:A2}
%     \EE\!\left[\calL_\rho^{k+1}\right]
%     \leq \EE\!\left[\calL_\rho^k\right]
%     - C_1^k\,\EE\!\left[\|\bx^{k+1}-\bx^k\|^2\right]
%     + \frac{\eta_k}{2}\,\EE\!\left[\|\be^k\|^2\right]
%     + \hat{\delta}_k,
%     \tag{A}
% \end{equation}
% where $C_1^k := \frac{1}{2\eta_k} - \frac{L_\varphi}{2} - \rho L_c^2 > 0$
% by~\eqref{eq:step-cond}.

% \medskip\noindent\textbf{Step 3: Scaled STORM recursion.}
% Multiply Lemma~\ref{lem:storm}~\eqref{eq:storm-error-bound} by
% $\frac{1}{8\eta_1 L_f^2}$:
% \begin{equation}\label{eq:B2}
%     \frac{\EE[\|\be^{k+1}\|^2]}{8\eta_1 L_f^2}
%     \leq
%     \frac{(1-\gamma_k)^2}{8\eta_1 L_f^2}\EE\!\left[\|\be^k\|^2\right]
%     + \frac{(1-\gamma_k)^2}{4\eta_1}\EE\!\left[\|\bx^{k+1}-\bx^k\|^2\right]
%     + \frac{\gamma_k^2\sigma^2}{4\eta_1 L_f^2}.
%     \tag{B}
% \end{equation}

% \medskip\noindent\textbf{Step 4: Combine and identify coefficients.}
% Adding \eqref{eq:A2} and \eqref{eq:B2}:
% \begin{align}
%     \varphi^{k+1}
%     &\leq \varphi^k
%     - \underbrace{\Bigl(C_1^k - \frac{(1-\gamma_k)^2}{4\eta_1}\Bigr)}_{=:\,\tilde{C}_1^k\,\geq\, 0}
%       \EE\!\left[\|\bx^{k+1}-\bx^k\|^2\right] \notag\\
%     &\quad
%     - \underbrace{\Bigl(\frac{\gamma_k(2-\gamma_k)}{8\eta_1 L_f^2}
%       - \frac{\eta_k}{2}\Bigr)}_{=:\,\tilde{C}_2^k\,\geq\, 0}
%       \EE\!\left[\|\be^k\|^2\right]
%     + \hat{\delta}_k
%     + \frac{\gamma_k^2\sigma^2}{4\eta_1 L_f^2}.
%     \label{eq:one-step}
% \end{align}
% Here we used $1-(1-\gamma_k)^2 = \gamma_k(2-\gamma_k)$.
% Positivity of $\tilde{C}_2^k$ follows from the second condition
% in~\eqref{eq:step-cond}:
% $\eta_k \leq \frac{\sqrt{\gamma_k(2-\gamma_k)}}{2L_f}$ implies
% $\eta_k^2 \leq \frac{\gamma_k(2-\gamma_k)}{4L_f^2}$, hence
% $\frac{\eta_k}{2} \leq \frac{\gamma_k(2-\gamma_k)}{8\eta_k L_f^2}
% \leq \frac{\gamma_k(2-\gamma_k)}{8\eta_1 L_f^2}$,
% where the last inequality uses $\eta_1 \geq \eta_k$.

% \medskip\noindent\textbf{Step 5: Drop non-negative terms and rearrange.}
% Since $\tilde{C}_1^k \geq 0$, dropping the $\|\bx^{k+1}-\bx^k\|^2$ term and
% rearranging~\eqref{eq:one-step}:
% \begin{equation}\label{eq:rearranged}
%     \tilde{C}_2^k\,\EE\!\left[\|\be^k\|^2\right]
%     \leq \varphi^k - \varphi^{k+1} + \hat{\delta}_k
%     + \frac{\gamma_k^2\sigma^2}{4\eta_1 L_f^2}.
% \end{equation}
% Since $\tilde{C}_2^k \geq \frac{\gamma_k(2-\gamma_k)}{8\eta_1 L_f^2}
% - \frac{\eta_k}{2} \geq \frac{\gamma_k(2-\gamma_k)}{16\eta_1 L_f^2}$
% (using $\frac{\eta_k}{2} \leq \frac{\gamma_k(2-\gamma_k)}{16\eta_1 L_f^2}$
% which follows from the step size condition), and since
% $\gamma_k(2-\gamma_k) \geq \gamma_k$ for $\gamma_k \in (0,1)$, we have
% $\tilde{C}_2^k \geq \frac{\gamma_k}{16\eta_1 L_f^2}$.
% Substituting into~\eqref{eq:rearranged}:
% \begin{equation}
%     \frac{\gamma_k}{16\eta_1 L_f^2}\,\EE\!\left[\|\be^k\|^2\right]
%     \leq \varphi^k - \varphi^{k+1} + \hat{\delta}_k
%     + \frac{\gamma_k^2\sigma^2}{4\eta_1 L_f^2}.
% \end{equation}

% \medskip\noindent\textbf{Step 6: Telescoping.}
% Summing over $k = 0, \ldots, K-1$:
% \begin{equation}
%     \frac{1}{16\eta_1 L_f^2}\sum_{k=0}^{K-1}\gamma_k\,\EE\!\left[\|\be^k\|^2\right]
%     \leq \varphi^0 - \varphi^K
%     + \sum_{k=0}^{K-1}\hat{\delta}_k
%     + \frac{1}{4\eta_1 L_f^2}\sum_{k=0}^{K-1}\gamma_k^2\sigma^2.
% \end{equation}
% Using $\varphi^K \geq \calL_\rho^* + \frac{\EE[\|\be^K\|^2]}{8\eta_1 L_f^2} \geq \calL_\rho^*$
% and expanding $\varphi^0 = \calL_\rho^0 + \frac{\EE[\|\be^0\|^2]}{8\eta_1 L_f^2}$
% where $\EE[\|\be^0\|^2] \leq 4\sigma^2$ (bounded by initialization), then
% multiplying both sides by $\frac{16\eta_1 L_f^2}{K}$ gives
% \begin{equation}
%     \frac{1}{K}\sum_{k=0}^{K-1}\gamma_k\,\EE\!\left[\|\be^k\|^2\right]
%     \leq
%     \frac{16\eta_1 L_f^2(\calL_\rho^0 - \calL_\rho^*)}{K}
%     + \frac{2\,\EE[\|\be^1\|^2]}{K}
%     + \frac{16\eta_1 L_f^2}{K}\sum_{k=0}^{K-1}\hat{\delta}_k
%     + \frac{1}{K}\sum_{k=0}^{K-1}4\gamma_k^2\sigma^2,
% \end{equation}
% which is~\eqref{eq:main-bound}.
% \end{proof}
\begin{lemma}[KKT residual stationarity]
\begin{align}
&\frac{1}{K}\sum_{k=0}^{K-1}\mathbb{E}\!\left[\|\zeta_x^{k+1}\|^2\right] \notag \\
\leq &
\left(\frac{\eta_1 L_f^2}{\gamma_K \cdot K} + \frac{1}{\eta_k \cdot K}\right)
\!\bigl(\mathcal{L}_\rho^0 - \mathcal{L}^*\bigr)
+ \frac{2\,\mathbb{E}\!\left[\|\mathbf{e}^1\|^2\right]}{\gamma_K \cdot K}
+ \frac{\displaystyle\sum_{k=0}^{K-1} 4\gamma_k^2\sigma^2}{\gamma_K \cdot K}
+ \frac{1}{K}\sum_{k=0}^{K-1}\delta_k
+ \left[\frac{1}{\eta_K \cdot K} + \frac{16\eta_1 L_f^2}{\gamma_K \cdot K}\right]
\sum_{k=0}^{K-1}\hat{\delta}_k
\end{align}
    
\end{lemma}
\begin{lemma}[Violation Bound]\label{lem:feasibility}
Under the same conditions as Lemma~\ref{lem:sufficient-descent},
with $\delta_k = c/(k+1)$ for some $c > 0$, we have
\begin{equation}\label{eq:feasibility-rate}
    \frac{1}{T}\sum_{k=0}^{T-1}\mathbb{E}\!\left[
        \|\bc_{\calE}(\bx^{k+1})\|^2
        + \|[\bc_{\calI}(\bx^{k+1})]_+\|^2
    \right]
    = \mathcal{O}\!\left(\frac{\log T}{T}\right)
\end{equation}
\end{lemma}

\begin{proof}
Let $r_k := \blmbd^k - \bmu^k$. By the $\blmbd$-update
\eqref{eq:lambda-update}, the equality block satisfies
$r_{\calE}^{k+1} = \rho\,\bc_{\calE}(\bx^{k+1})$, while the inequality block
obeys
\[
    r_{\calI}^{k+1}
    = \bigl[\bmu_{\calI}^{k+1} + \rho\,\bc_{\calI}(\bx^{k+1})\bigr]_+
    - \bmu_{\calI}^{k+1}.
\]
Since $\bmu_{\calI}^{k+1}\geq \bm{0}$ by the initialization and the convex
combination update \eqref{eq:mu-update}, we have
$[c_i(\bx^{k+1})]_+ \leq |r_i^{k+1}|/\rho$ for every $i\in\calI$, and hence
\begin{equation}\label{eq:r-feasibility}
  \|\bc_{\calE}(\bx^{k+1})\|^2 + \|[\bc_{\calI}(\bx^{k+1})]_+\|^2
  \leq \frac{1}{\rho^2}\|r_{k+1}\|^2,
\end{equation}
so it suffices to bound $\frac{1}{T}\sum_{k=0}^{T-1}\mathbb{E}[\|r_{k+1}\|^2]$.

\medskip\noindent\textbf{Step 1: Derive the recursion for $\|r_{k+1}\|^2$.}
By the triangle inequality and $(a+b)^2 \leq 2(a^2+b^2)$,
\begin{equation}\label{eq:r-triangle}
  \|r_{k+1}\|^2
  \leq 2\|\blmbd^{k+1} - \blmbd^k\|^2 + 2\|\blmbd^k - \bmu^{k+1}\|^2.
\end{equation}
Applying Lemma~\ref{lem:basic-relations}~\eqref{eq:mu-lambda-diff}
to the second term and \eqref{eq:lambda-diff} to the first:
\begin{align}
  \|\blmbd^k - \bmu^{k+1}\|^2 &= (1-\theta_k)^2\|r_k\|^2, \notag\\
  \mathbb{E}[\|\blmbd^{k+1}-\blmbd^k\|^2 \mid \mathcal{F}_k]
  &\leq 2\rho^2 L_c^2\,\mathbb{E}[\|\bx^{k+1}-\bx^k\|^2 \mid \mathcal{F}_k]
    + \frac{\rho^2\delta_k^2}{2}. \notag
\end{align}
Substituting into \eqref{eq:r-triangle} and using
$(1-\theta_k)^2 \leq (1-\theta_k)$ (since $\theta_k \in (0,1]$):
\begin{equation}\label{eq:r-recursion}
  \mathbb{E}[\|r_{k+1}\|^2 \mid \mathcal{F}_k]
  \leq 2(1-\theta_k)\|r_k\|^2
       + 4\rho^2 L_c^2\,\mathbb{E}[\|\bx^{k+1}-\bx^k\|^2 \mid \mathcal{F}_k]
       + \rho^2\delta_k^2.
\end{equation}

\medskip\noindent\textbf{Step 2: Sum the recursion.}
Rearranging \eqref{eq:r-recursion}:
\begin{equation}
  2\theta_k\|r_k\|^2
  \leq 2\|r_k\|^2 - \mathbb{E}[\|r_{k+1}\|^2 \mid \mathcal{F}_k]
       + 4\rho^2 L_c^2\,\mathbb{E}[\|\bx^{k+1}-\bx^k\|^2 \mid \mathcal{F}_k]
       + \rho^2\delta_k^2.
\end{equation}
Taking total expectation and summing over $k = 0,\ldots,T-1$:
\begin{equation}\label{eq:theta-r-sum}
  2\sum_{k=0}^{T-1}\theta_k\,\mathbb{E}[\|r_k\|^2]
  \leq 2\sum_{k=0}^{T-1}\mathbb{E}[\|r_k\|^2]
       - \sum_{k=0}^{T-1}\mathbb{E}[\|r_{k+1}\|^2]
       + 4\rho^2 L_c^2\sum_{k=0}^{T-1}\mathbb{E}[\|\bx^{k+1}-\bx^k\|^2]
       + \rho^2\sum_{k=0}^{T-1}\delta_k^2.
\end{equation}
Note that the first two sums on the right telescope:
\begin{equation}\label{eq:telescope-r}
  2\sum_{k=0}^{T-1}\mathbb{E}[\|r_k\|^2] - \sum_{k=0}^{T-1}\mathbb{E}[\|r_{k+1}\|^2]
  = \sum_{k=0}^{T-1}\mathbb{E}[\|r_k\|^2] - \sum_{k=1}^{T-1}\mathbb{E}[\|r_k\|^2]
    + \mathbb{E}[\|r_0\|^2] - \mathbb{E}[\|r_T\|^2]
  \leq \mathbb{E}[\|r_0\|^2].
\end{equation}
Substituting \eqref{eq:telescope-r} into \eqref{eq:theta-r-sum}:
\begin{equation}\label{eq:theta-r-final}
  2\sum_{k=0}^{T-1}\theta_k\,\mathbb{E}[\|r_k\|^2]
  \leq \|r_0\|^2
       + 4\rho^2 L_c^2\sum_{k=0}^{T-1}\mathbb{E}[\|\bx^{k+1}-\bx^k\|^2]
       + \rho^2\sum_{k=0}^{T-1}\delta_k^2.
\end{equation}
Dividing both sides of \eqref{eq:theta-r-final} by $T$:
  \textcolor{red}{
\begin{equation}\label{eq:feasibility-weighted}
  \frac{2}{T}\sum_{k=0}^{T-1}\theta_k\,\mathbb{E}[\|r_k\|^2]
  \leq \frac{\|r_0\|^2}{T}
       + \frac{4\rho^2 L_c^2}{T}\sum_{k=0}^{T-1}\mathbb{E}[\|\bx^{k+1}-\bx^k\|^2]
       + \frac{\rho^2}{T}\sum_{k=0}^{T-1}\delta_k^2.
\end{equation}}
% where the last equality follows from
% Lemma~\ref{lem:sufficient-descent} (which gives
% $\frac{1}{T}\sum_{k=0}^{T-1}\mathbb{E}[\|\bx^{k+1}-\bx^k\|^2]
% = \mathcal{O}(\log T / T)$)
% and $\frac{1}{T}\sum_{k=0}^{T-1}\delta_k^2 = \mathcal{O}(1/T)$
% for $\delta_k = c/(k+1)$.
% Since $\theta_k \geq 0$ and $\|r_k\|^2 = \rho^2\|\bc(\bx^k)\|^2 \geq 0$,
% \eqref{eq:feasibility-weighted} implies that
% $\theta_k\,\mathbb{E}[\|\bc(\bx^k)\|^2] \to 0$ in a weighted ergodic sense.

% \medskip\noindent\textbf{Step 4: Bound each term.}
% By Lemma~\ref{lem:sufficient-descent} and Lemma~\ref{lem:storm}:
% \begin{equation}
%   \frac{1}{T}\sum_{k=0}^{T-1}\mathbb{E}[\|\bx^{k+1}-\bx^k\|^2]
%   = \mathcal{O}\!\left(\frac{\log T}{T}\right).
% \end{equation}
% For $\delta_k = c/(k+1)$:
% \[
%   \frac{1}{T}\sum_{k=0}^{T-1}\delta_k = \mathcal{O}\!\left(\frac{\log T}{T}\right),
%   \qquad
%   \frac{1}{T}\sum_{k=0}^{T-1}\delta_k^2 = \mathcal{O}\!\left(\frac{1}{T}\right).
% \]
% Substituting into \eqref{eq:feasibility-pre} and applying \eqref{eq:r-feasibility}
% yields \eqref{eq:feasibility-rate}.
\end{proof}

\section{Numerical Experiments}
\label{sec:experiments}

\appendix
\section{Supported Lemmas}
\label{sec:appendix-supported}

\begin{lemma}[Vector Azuma-Hoeffding Inequality]
\label{lem:vector-azuma}
Let $\{\bv^k\}_{k=1}^K$ be a vector-valued martingale difference
sequence adapted to filtration $\{\calF_k\}_{k=0}^K$, satisfying
\begin{equation}
    \EE[\bv^k \mid \calF_{k-1}] = \zero,
    \qquad
    \|\bv^k\| \leq a_k,\quad k=1,\ldots,K.
\end{equation}
Then for any $\delta \in (0,1)$, with probability at least $1-\delta$,
\begin{equation}
    \left\|\sum_{k=1}^{K}\bv^k\right\|
    \leq
    3\sqrt{\log(1/\delta)\sum_{k=1}^{K}a_k^2}.
\end{equation}
\end{lemma}

\noindent
The above bound is a standard vector-valued Azuma-Hoeffding concentration result.

\clearpage
\small
\bibliographystyle{plain}
\bibliography{refer}
\end{document}
